\section{Conclusion}
\label{sec:conclusion}

We began by observing that combining formal verification strategies for LLM code generation \emph{appears} straightforward---stack grammar constraints, static analysis, and SMT-guided repair for multiplicative gains. Our work validates this intuition in part, while revealing a critical prerequisite that prior work overlooked.

\subsection{Summary}

In this work, we addressed whether verification strategies target independent error classes, enabling effective pipeline combination. Our main contributions are:

\begin{enumerate}
    \item \textbf{Error-class-independence framework.} We provided the first quantitative measurement showing that grammar, semantic, and specification error classes are largely independent (mean Jaccard 0.0752, all pairs below 0.30). This validates the theoretical foundation for layered verification.

    \item \textbf{Model capability prerequisite.} We identified that feedback-based verification requires instruction-tuned models. A completion model given static analysis feedback increased security issues by 600\%, demonstrating that detection is model-agnostic while repair is model-capability-dependent.

    \item \textbf{Partial pipeline validation.} We demonstrated that grammar-constrained decoding reduces syntax errors by 40\%, validating the first stage of the verification pipeline.
\end{enumerate}

\subsection{Future Directions}

Our experiments point to several promising directions grounded in specific findings:

\textbf{From H-M2 failure:} The static analysis feedback loop failed with StarCoder2-3b. Future work should retest with instruction-tuned models (Llama-3-8B-Instruct, CodeLlama-7B-Instruct) to determine whether model capability alone explains the failure, or whether prompt formatting also plays a role.

\textbf{From unverified assumptions:} We tested only the syntax$\rightarrow$semantics$\rightarrow$specifications ordering. Alternative orderings (parallel application, different sequences) remain untested and might yield better results for specific error distributions.

\textbf{From scope limitations:} All experiments used single-function Python generation. Extending to repository-level code generation, multi-file projects, and additional programming languages would test generalization of the independence framework.

\textbf{Full pipeline validation:} With instruction-tuned models, completing H-M3 (SMT repair) and H-M4 (synergy measurement) would enable computing the synergy coefficient and testing the multiplicative improvement hypothesis directly.

\subsection{Closing}

Layered verification for LLM code generation is viable---error classes are genuinely independent, and each verification stage contributes non-redundant value. But model capability is the hidden gate: grammar constraints work during generation, while feedback loops require models that understand repair instructions. We hope this distinction helps practitioners deploy verification pipelines that improve code quality rather than inadvertently degrading it.
